\documentclass[12pt]{article}

\usepackage[margin=2.5cm]{geometry}
\usepackage{amsmath, amssymb}
\usepackage{graphicx}
\usepackage{booktabs}
\usepackage{hyperref}
\usepackage[numbers]{natbib}
\usepackage{lineno}
\usepackage{setspace}
\usepackage{xcolor}

\doublespacing
\linenumbers

\hypersetup{
  colorlinks=true,
  linkcolor=blue,
  citecolor=blue,
  urlcolor=blue
}

% ─────────────────────────────────────────
\title{\textbf{Aging as Synchronization Failure: A Minimal Model\\
Explaining the Scaling Limits of Anti-Aging Interventions}}

\author{
  Piotr Czarnecki\\
  \small Independent Researcher\\
  \small \texttt{piotrczary@gmail.com}
}

\date{March 2026 \\ \vspace{0.5em}
  \small\textit{Preprint — not yet peer reviewed}}

% ─────────────────────────────────────────
\begin{document}
\maketitle

% ─────────────────────────────────────────
\begin{abstract}
We present a minimal computational model of biological aging based on the
dynamics of coupled oscillators. The model proposes that aging emerges primarily
from progressive desynchronization among biological regulatory systems, rather
than from the accumulation of component-level damage alone. A positive feedback
loop---increasing variance $\sigma$ raises synchronization cost, which depletes
the regulatory buffer $B$ and coupling strength $K$, which further increases
$\sigma$---produces aging dynamics consistent with empirically observed patterns
including the Gompertz mortality law. Crucially, the model generates a
quantitative prediction regarding intervention efficacy: local repair of a single
oscillator reduces system-wide desynchronization by approximately
$\Delta\sigma/\sigma \sim (2f-f^2)/n$, where $f$ is repair effectiveness and $n$ is
system complexity. This scaling law implies that local anti-aging interventions,
highly effective in simple systems such as mice, lose therapeutic significance in
complex systems such as humans. Global synchronization interventions, by
contrast, maintain constant effectiveness independent of $n$. The model provides
a mechanistic account of the well-documented failure to translate anti-aging
findings from animal models to humans.
\end{abstract}

\textbf{Keywords:} aging, Kuramoto oscillators, synchronization, Gompertz law,
anti-aging interventions, systems biology

\newpage

% ─────────────────────────────────────────
\section{Introduction}

The dominant paradigm in aging research focuses on the accumulation of molecular
damage: telomere shortening, mitochondrial dysfunction, DNA error accumulation,
and oxidative stress \citep{lopez2023hallmarks}. This framework has produced
valuable insights but faces a persistent explanatory gap: damage accumulates
approximately linearly, yet aging accelerates nonlinearly. An eighty-year-old is
not simply ``twice as old'' as a forty-year-old in any functional sense---they
occupy a qualitatively different physiological regime.

A complementary perspective focuses not on damage per se but on the loss of
coordination among biological systems. With age, circadian rhythms desynchronize
from metabolic rhythms; cortisol cycles lose precision; neural networks that
operate coherently in youth become progressively decoupled. The variance of
biological parameters---glucose levels, blood pressure, hormone
concentrations---increases with age not because mean values deteriorate
dramatically, but because the systems that regulate these parameters lose mutual
coherence.

This observation motivates a synchronization-theoretic approach to aging. The
Kuramoto model, originally developed to describe the spontaneous synchronization
of coupled oscillators \citep{kuramoto1975}, provides a natural mathematical
framework. In this framework, the biological organism is modeled as a network of
oscillators---each representing a major regulatory system---that maintain
coherence through coupling forces. Aging corresponds to the progressive loss of
this coherence, driven by a positive feedback loop between desynchronization,
energetic cost, and structural damage.

Independent empirical support for this view comes from longitudinal analysis of
blood biomarkers, which reveals that the time required for organisms to return to
homeostasis after perturbation increases approximately fourfold between ages 40
and 90 \citep{pyrkov2021}---a signature of systems approaching a critical
transition, consistent with the synchronization framework.

The present work introduces a minimal implementation of this framework, derives
its key quantitative predictions, and explores their implications for the design
of anti-aging interventions.

% ─────────────────────────────────────────
\section{Model}

\subsection{Biological oscillators}

Each organism is represented as a network of $N$ coupled oscillators with phases
$\theta_i(t)$ governed by the Kuramoto equation with noise
\citep{kuramoto1975, strogatz2000}:

\begin{equation}
\frac{d\theta_i}{dt} = \omega_i + \frac{K}{N} \sum_j \sin(\theta_j - \theta_i) + \eta_i(t)
\end{equation}

where $\omega_i \sim \text{Uniform}(0.9, 1.1)$ are natural frequencies, $K$ is
coupling strength, and $\eta_i(t)$ is Gaussian noise with standard deviation
$\sigma_\text{noise}$.

The degree of synchronization across the network is measured by the order
parameter $r = |\langle e^{i\theta} \rangle|$, and desynchronization is defined
as:

\begin{equation}
\sigma = 1 - r \in [0, 1]
\end{equation}

where $\sigma = 0$ indicates perfect synchronization and $\sigma = 1$ indicates
complete incoherence.

\subsection{Damage feedback loop}

The key innovation of the model is the introduction of a positive feedback loop
connecting desynchronization to structural damage:

\begin{align}
\text{Cost} &= \frac{\sigma^{3.5}}{B} \\[4pt]
\frac{dD}{dt} &= \text{Cost} \\[4pt]
\frac{dB}{dt} &= -\alpha \cdot \text{Cost} \\[4pt]
\frac{dK}{dt} &= -\beta \cdot \text{Cost} \cdot (1 + \sigma)
\end{align}

where $D$ is accumulated damage, $B$ is the regulatory buffer (initial value
1.0), and $\alpha$, $\beta$ are degradation parameters. The nonlinear cost
function ($\sigma^{3.5}$) ensures that small increases in desynchronization
produce disproportionately large energetic costs at high $\sigma$, consistent
with the Jensen's inequality argument for multiplicative biological systems.

Additionally, the effective noise amplitude grows exponentially with accumulated
damage:

\begin{equation}
\eta_\text{eff} = \eta_0 \cdot \min\!\left(\exp\!\left(\frac{0.25 \cdot D}{D_\text{scale}}\right),\; 5.0\right)
\end{equation}

This captures the empirical observation that aging systems become progressively
less able to buffer perturbations, producing increasing variability in biological
parameters \citep{pyrkov2021}.

\subsection{Organism parameters}

Two canonical organisms are implemented with the parameters shown in
Table~\ref{tab:params}.

\begin{table}[h]
\centering
\caption{Model parameters for mouse and human.}
\label{tab:params}
\begin{tabular}{lcc}
\toprule
Parameter & Mouse ($n{=}5$) & Human ($n{=}20$) \\
\midrule
Oscillators $N$ & 5 & 20 \\
Coupling $K_0$ & 0.80 & 0.30 \\
Noise $\sigma_0$ & 0.05 & 0.15 \\
$\alpha$ & 0.0015 & 0.0020 \\
$\beta$ & 0.003 & 0.004 \\
\bottomrule
\end{tabular}
\end{table}

We emphasize that these mappings are illustrative rather than calibrated: the
distinction between $n{=}5$ and $n{=}20$ represents a qualitative difference in
regulatory complexity, not a literal count of independent biological oscillators.
An operational measure of $n$ can be defined empirically as the number of
principal components explaining $\geq\!90\%$ of variance in a panel of
biological markers — a quantity directly measurable in existing biobank datasets
such as UK Biobank or NHANES, and one that would allow the scaling law to be
tested quantitatively against real intervention data.

% ─────────────────────────────────────────
\section{Results}

\subsection{Emergence of the Gompertz mortality law}

The qualitative behavior of the model is robust across a wide range of parameter
values (noise amplitude, coupling strength, $\alpha$, $\beta$): both the
emergence of Gompertz-like mortality dynamics and the $1/n$ scaling law are
preserved. Furthermore, both the nonlinear cost function ($\sigma^{3.5}$) and
the exponential noise feedback are necessary for reproducing Gompertz-like
dynamics; removing either eliminates the observed acceleration of mortality risk
and produces approximately linear damage accumulation instead.

Simulation of cohorts of 200 agents with death defined as $B < 0.02$ or $\sigma
> 0.95$ produces survival curves consistent with empirically observed mortality
patterns. Critically, the log-hazard $h(t)$ follows an approximately linear
trend in time---the defining characteristic of the Gompertz mortality
law---despite this law not being explicitly encoded in the model. This emergence
occurs because the exponential noise feedback creates exponentially accelerating
desynchronization near the death threshold, producing the characteristic doubling
of mortality risk with age.

The median survival time differs between mouse ($t_{50} \approx 80$) and human
($t_{50} \approx 72$) in model time units. These units are dimensionless and
cannot be directly compared to biological lifespan in years; the model captures
the \emph{relative} difference in aging trajectory between organisms of different
complexity, not absolute longevity. The human curve shows greater
late-life acceleration---consistent with the more complex coupling architecture
being less recoverable once desynchronization initiates.

\subsection{Complexity determines aging trajectory}

Figure~1 constitutes the central empirical result of the model. It shows the
desynchronization trajectory $\sigma(t)$ averaged over 15 independent
simulations for networks of varying size ($n = 3, 5, 8, 12, 20, 30$). The key
finding is monotonic: larger networks enter the high-desynchronization regime
earlier and more reliably. The standard deviation of $\sigma$ decreases with
$n$, indicating that complex systems, while more susceptible to
desynchronization, are more predictable in their aging trajectory.

This result emerges from the interplay between coupling architecture and noise
accumulation. In small networks, strong relative coupling can periodically
restore coherence, producing observed returns to low $\sigma$. In large networks,
the same absolute coupling represents a smaller fraction of the total
synchronization demand, making recovery progressively less likely.

\subsection{Scaling law for intervention efficacy}

The central quantitative prediction of the model concerns the effectiveness of
interventions that repair a single oscillator versus interventions that enhance
global coupling. We derive this law from first principles.

\subsubsection*{Derivation}

In the mean-field Kuramoto model with noise, each oscillator $i$ has phase
$\theta_i$. Near a synchronized state, phases fluctuate around the mean-field
phase $\Psi$ with small deviations $\phi_i = \theta_i - \Psi$. Expanding the
complex order parameter $Z = \frac{1}{N}\sum_j e^{i\theta_j}$ to second order:

\begin{equation}
r = |Z| \approx 1 - \frac{1}{2N}\sum_{j=1}^{N} \langle \phi_j^2 \rangle
\end{equation}

In the linearized dynamics, the steady-state phase variance of oscillator $j$ is
proportional to its squared noise amplitude:
$\langle \phi_j^2 \rangle = c \cdot \eta_j^2$,
where $c$ is a constant depending on coupling strength and frequency
distribution. Therefore:

\begin{equation}
\sigma = 1 - r \approx \frac{c}{2N} \sum_{j=1}^{N} \eta_j^2
\end{equation}

\textit{Proposition.} Let all oscillators have identical noise $\eta_0$, so that
$\sigma_0 = \frac{c \eta_0^2}{2}$. Suppose a local intervention reduces the noise
of a single oscillator $k$ by fraction $f \in [0,1]$, setting
$\eta_k \leftarrow (1-f)\eta_0$. Then:

\begin{equation}
\boxed{\frac{\Delta\sigma}{\sigma_0} = \frac{f(2-f)}{N}} \quad \text{(local)}
\end{equation}

\textit{Proof.} After the intervention:
\begin{align}
\sigma_\text{new} &= \frac{c}{2N}\left[(N-1)\eta_0^2 + (1-f)^2\eta_0^2\right] \notag\\
&= \frac{c\eta_0^2}{2N}\left[N - 2f + f^2\right]
\end{align}

Therefore:
\begin{align}
\Delta\sigma &= \sigma_0 - \sigma_\text{new}
= \frac{c\eta_0^2}{2N}\left[N - (N - 2f + f^2)\right]
= \frac{c\eta_0^2}{2N}(2f - f^2) \notag
\end{align}

\begin{equation}
\frac{\Delta\sigma}{\sigma_0}
= \frac{c\eta_0^2(2f-f^2)/(2N)}{c\eta_0^2/2}
= \frac{f(2-f)}{N} \qquad \square
\end{equation}

\textit{Global intervention.} If the same fractional repair $f$ is applied to
all $N$ oscillators simultaneously:
\begin{equation}
\sigma_\text{new}^\text{global} = \frac{c\eta_0^2}{2}(1-f)^2
\end{equation}
\begin{equation}
\frac{\Delta\sigma}{\sigma_0}\bigg|_\text{global} = 1 - (1-f)^2 = f(2-f)
\end{equation}

This is independent of $N$. The ratio of local to global efficacy is therefore
exactly $1/N$:

\begin{equation}
\frac{(\Delta\sigma/\sigma_0)_\text{local}}{(\Delta\sigma/\sigma_0)_\text{global}} = \frac{1}{N}
\end{equation}

We note a correction to our earlier shorthand: the precise formula is
$f(2-f)/N$, not $f^2/N$. The two expressions agree for $f \ll 1$ (small repair)
and for $f=1$ (complete repair) give $1/N$ vs $1/N$, but differ at intermediate
$f$. Table~\ref{tab:scaling} reports the exact values.

\subsubsection*{Numerical values}

For $f = 0.5$ (50\% noise reduction), $f(2-f) = 0.75$:

\begin{table}[h]
\centering
\caption{Exact predicted reduction $\Delta\sigma/\sigma$ for local vs.\
global interventions ($f = 0.5$, $f(2-f) = 0.75$).}
\label{tab:scaling}
\begin{tabular}{lcc}
\toprule
System ($n$) & Local: $\frac{f(2-f)}{n}$ & Global: $f(2-f)$ \\
\midrule
3 (simple organism) & 25.0\% & 75.0\% \\
5 (mouse)           & 15.0\% & 75.0\% \\
12                  & 6.3\%  & 75.0\% \\
20 (human)          & 3.8\%  & 75.0\% \\
30                  & 2.5\%  & 75.0\% \\
\bottomrule
\end{tabular}
\end{table}

The global intervention maintains constant efficacy regardless of system
complexity. This divergence between local and global efficacy, scaling exactly
as $1/N$, constitutes the primary falsifiable prediction of the model.

\subsubsection*{Robustness of the scaling}

The $1/N$ ratio between local and global efficacy is exact within the
mean-field approximation and does not depend on the specific values of coupling
$K$, noise $\eta_0$, frequency distribution, or the repair fraction $f$. It
holds as long as (i) noise contributions to the order parameter are additive,
and (ii) oscillators are statistically exchangeable. The first assumption holds
in the linearized regime; the second is violated in hierarchical networks
(a limitation discussed in Section~4.3).

We verified numerically that the qualitative scaling is preserved when the cost
exponent is varied between 2 and 5, when $\alpha$ and $\beta$ are varied by
factors of 2--5, and when $n$ varies continuously from 3 to 30.

% ─────────────────────────────────────────
\section{Discussion}

\subsection{The translational gap in anti-aging research}

The failure to translate anti-aging interventions from animal models to humans is
one of the most consistent and costly problems in aging research. Interventions
targeting specific molecular pathways---caloric restriction mimetics, senolytics,
NAD\textsuperscript{+} precursors, mTOR inhibitors---show dramatic effects in
short-lived organisms but modest or negligible effects in clinical trials
\citep{barzilai2018}.

The scaling law derived here provides a mechanistic explanation for this failure
that requires no assumption of fundamental biological difference between species.
A local intervention is inherently less effective in a more complex organism, not
because the target pathway is less important, but because any single pathway
represents a smaller fraction of the system's total synchronization demand. An
intervention that addresses 25\% of synchronization variance in a
three-oscillator system ($n{=}3$, $f{=}0.5$: $f(2{-}f)/n = 0.75/3 = 25\%$)
addresses only 3.8\% in a twenty-oscillator system ($n{=}20$: $0.75/20 = 3.8\%$),
at
identical molecular effectiveness.

\subsection{Implications for intervention design}

The model suggests a reorientation of anti-aging strategy from local repair
toward global synchronization. Interventions that strengthen the coupling among
regulatory systems---rather than repairing any specific system---maintain
efficacy independent of organism complexity. Empirically, this corresponds to
interventions including regular physical exercise at consistent times, structured
sleep protocols, intermittent fasting, and other zeitgeber-based approaches.
These interventions are notable for producing consistent (if modest) benefits
across both animal models and human clinical populations \citep{rakshit2023},
consistent with complexity-independent efficacy.

This reorientation does not imply that local molecular interventions are without
value. It implies that their value is primarily diagnostic and mechanistic rather
than therapeutic---they identify the components of biological synchronization
networks, but the therapeutic target is the network itself.

\subsection{Model limitations}

The present model is deliberately minimal and makes several simplifying
assumptions that warrant explicit acknowledgment.

First, all oscillators are treated as equivalent in role and coupling
architecture. Biological regulatory networks have hierarchical structure---the
suprachiasmatic nucleus serves as a master oscillator---and this hierarchy is not
captured in the current implementation.

Second, the mapping between model parameters and biological quantities is
illustrative rather than calibrated.

Third, the model does not account for developmental processes, genetic variation,
or environmental heterogeneity that substantially influence aging in both animal
models and human populations.

These limitations suggest directions for model extension rather than fundamental
challenges to the core framework. The model should be interpreted as a minimal
(toy) model capturing qualitative system-level dynamics, rather than a
quantitatively calibrated physiological simulator.

\subsection{Testable predictions}

The model generates three predictions testable with existing methods:

\textbf{Prediction 1 (primary):} The effectiveness of single-pathway
interventions in reducing biological age markers scales approximately as $1/n$,
where $n$ is a measure of regulatory network complexity.

\textbf{Prediction 2:} The cross-correlational structure of biological
parameters---the degree to which circadian, metabolic, immune, and hormonal
rhythms are mutually correlated---should be a better predictor of functional age
than the mean values of any single parameter.

\textbf{Prediction 3:} Zeitgeber-based interventions should show greater
relative benefit in older organisms with high existing desynchronization than in
younger organisms with intact synchronization, because the marginal value of
coupling restoration is highest near the desynchronization threshold.

% ─────────────────────────────────────────
\section{Conclusion}

We have presented a minimal model in which aging emerges as a positive feedback
loop of desynchronization in coupled biological oscillators. The model reproduces
the Gompertz mortality law without encoding it directly, generates differential
aging trajectories as a function of system complexity, and produces the
quantitative scaling law $\Delta\sigma/\sigma = (2f-f^2)/n$ for local repair
interventions.

The primary implication is practical: the consistent failure to translate
anti-aging interventions from simple organisms to humans is not a biological
mystery requiring species-specific explanation---it is a mathematical consequence
of the scaling properties of synchronization dynamics. Interventions targeting
specific components of a synchronization network necessarily lose efficacy as
network complexity grows. Interventions targeting the network coupling itself do
not.

Aging, in this framework, is not primarily a story of damage accumulation. It is
a story of a network gradually losing its ability to coordinate. The most direct
therapeutic target is the coordination itself.

This suggests that aging interventions should be evaluated not only by their
local biochemical effects, but by their impact on system-wide synchronization---a
criterion largely absent from current clinical trial design.

% ─────────────────────────────────────────
\section*{Code availability}

Simulation code is available at:
\url{https://github.com/piotrczary/aging-synchronization}

\section*{Competing interests}

The author declares no competing interests.

\appendix

% ─────────────────────────────────────────
\section{Scaling of Interventions in the Kuramoto Framework}
\label{app:scaling}

\subsection{Variance approximation of desynchronization}

Consider a system of $n$ coupled oscillators with phases $\theta_i$. The
Kuramoto order parameter is:
\begin{equation}
r = \left|\frac{1}{n}\sum_{i=1}^{n} e^{i\theta_i}\right|
\end{equation}

We define macroscopic desynchronization as $\sigma = 1 - r$. For a synchronized
state with small phase dispersion, Taylor expansion yields
$e^{i\theta_i} \approx 1 + i\theta_i - \theta_i^2/2$. Assuming the mean phase
is centered ($\langle\theta_i\rangle = 0$), the real part of the order
parameter is approximately:
\begin{equation}
r \approx 1 - \frac{1}{2}\,\mathrm{Var}(\theta)
\end{equation}

Thus, system desynchronization scales directly with phase variance:
$\sigma \approx \frac{1}{2}\,\mathrm{Var}(\theta)$.

\subsection{Noise-driven variance decomposition}

Assuming each oscillator is perturbed by independent noise $\xi_i$ with
amplitude (standard deviation) $\eta_i$, the steady-state phase variance scales
as $\langle\phi_i^2\rangle = c\,\eta_i^2$, where $c$ depends on coupling
strength and frequency distribution. The total system desynchronization is
therefore:
\begin{equation}
\sigma = \frac{c}{2n}\sum_{i=1}^{n}\eta_i^2
\end{equation}

For uniform baseline noise $\eta_i = \eta_0$:
$\sigma_0 = \frac{c\,\eta_0^2}{2}$.

\subsection{Local vs.\ global interventions}

Let an intervention reduce the noise amplitude of target oscillator $k$ by
fraction $f \in (0,1]$, so $\eta_k \to (1-f)\eta_k$.

The change in variance contribution of oscillator $k$ is:
\begin{equation}
\Delta(\eta_k^2) = \eta_0^2 - (1-f)^2\eta_0^2 = \eta_0^2(2f - f^2)
\end{equation}

For a \textbf{local intervention} affecting only one oscillator out of $n$:
\begin{equation}
\frac{\Delta\sigma}{\sigma_0}
= \frac{c\,\eta_0^2(2f-f^2)/(2n)}{c\,\eta_0^2/2}
= \frac{2f - f^2}{n}
\label{eq:app_local}
\end{equation}

For a \textbf{global intervention} where all $n$ oscillators receive the same
fractional improvement $f$:
\begin{equation}
\frac{\Delta\sigma}{\sigma_0} = 1 - (1-f)^2 = 2f - f^2
\label{eq:app_global}
\end{equation}

This is independent of $n$.

\subsection{Conclusion}

The ratio of local to global efficacy is exactly $1/n$, independent of $f$,
$\eta_0$, coupling strength, or frequency distribution:
\begin{equation}
\frac{(\Delta\sigma/\sigma_0)_\text{local}}{(\Delta\sigma/\sigma_0)_\text{global}}
= \frac{1}{n}
\end{equation}

For a complete local repair ($f = 1$): $(2f - f^2)|_{f=1} = 1$, so local
efficacy reduces precisely to $1/n$. Even a hypothetically perfect intervention
eliminating all noise in a single oscillator is bounded by network complexity.
This structural property holds for any $n \geq 1$ and any $f \in (0,1]$.

\subsection{Regime of validity}

The derivation above relies on two approximations that warrant explicit
acknowledgment.

First, the Taylor expansion $r \approx 1 - \frac{1}{2}\mathrm{Var}(\theta)$
holds only for \emph{small phase dispersion}. In the main simulations, $\sigma$
reaches values approaching 0.95 — near-complete desynchronization — where this
linearization is quantitatively inaccurate. The analytical derivation therefore
rigorously establishes the $1/n$ scaling in the weakly-synchronized regime.
Numerical simulations confirm that the scaling generalizes beyond this regime
(Section~3.3), but we cannot claim analytical proof in the high-desynchronization
limit.

Second, the variance decomposition treats oscillator phases as independent.
However, the Kuramoto coupling term explicitly introduces phase correlations.
The derivation ignores coupling-induced covariances, which is justified in the
limit $K \to 0$ or when phase fluctuations are small relative to $2\pi$.

In summary: the $1/n$ scaling law is proven exactly within the weak-coupling,
small-dispersion approximation. Numerical simulations confirm that the
qualitative result persists across the full dynamical range of the model,
including the high-$\sigma$ regime where system failure occurs.

% ─────────────────────────────────────────
\section{Robustness of Gompertz Dynamics to Cost Function Form}
\label{app:gompertz}

\subsection{Objective}

The baseline model employs a cost function $C(\sigma) \propto \sigma^x$ with
$x = 3.5$. This appendix demonstrates that Gompertz-like mortality dynamics
emerge for any $x > 1$, depending only on two structural properties: (i)
convexity of the cost function, and (ii) positive feedback between accumulated
damage and noise amplitude.

\subsection{Generalized feedback structure}

Consider the system:
\begin{align}
\frac{dD}{dt} &= \frac{\sigma^x}{B} \\
\frac{dB}{dt} &= -\alpha\,\frac{\sigma^x}{B} \\
\frac{dK}{dt} &= -\beta\,\frac{\sigma^x}{B}(1 + \sigma)
\end{align}
with exponential noise scaling $\eta_\text{eff} = \eta_0 \cdot \exp(\gamma D)$,
$\gamma > 0$. The system contains a closed positive feedback loop:
\begin{equation}
\sigma \uparrow \;\Rightarrow\; C \uparrow \;\Rightarrow\; D \uparrow
\;\Rightarrow\; \eta \uparrow \;\Rightarrow\; \sigma \uparrow
\end{equation}

\subsection{Late-stage approximation and finite-time singularity}

Near the end of life, $B$ is small and noise dominates coupling. We approximate
$dD/dt \sim \sigma^x$. This approximation assumes $\sigma$ is dominated by noise
rather than coupling ($K \approx 0$ regime), which is justified late in the
simulation when $K$ has been depleted by accumulated damage. Due to exponential noise amplification,
$\sigma(t) \sim \exp(\gamma D(t))$, giving:
\begin{equation}
\frac{dD}{dt} \sim e^{x\gamma D}
\end{equation}

Separating variables and integrating:
\begin{equation}
D(t) \sim -\frac{1}{x\gamma}\log(t_c - t)
\end{equation}

for some critical time $t_c$, showing $D \to \infty$ as $t \to t_c$ — a
finite-time singularity characteristic of systems undergoing catastrophic failure.

\subsection{Implication for mortality hazard}

Since $\sigma(t) \sim \exp(\gamma D(t))$ accelerates toward the failure
threshold, the instantaneous hazard in the pre-critical regime satisfies:
\begin{equation}
h(t) \propto \exp(\lambda t), \quad \lambda > 0
\end{equation}

This is the Gompertz law: $\log h(t)$ grows linearly in time. The rate $\lambda$
depends on $x$ and the degradation parameters, but the \emph{qualitative
structure} is preserved for all $x > 1$.

\subsection{Independence from exponent $x$}

The exponent $x$ influences the rate of acceleration but not the form of the
dynamics. For any $x > 1$: the cost is convex, the feedback loop is active,
damage accelerates super-exponentially, and hazard grows approximately
exponentially. Gompertz-like dynamics therefore emerge from the architecture of
the model, not from parameter tuning.

\textbf{Key result:} Gompertz-like dynamics emerge for all $x > 1$.

\subsection{Numerical verification}

We simulated the model with cost exponents $x \in \{2.0, 2.5, 3.0, 3.5, 4.0,
5.0\}$. In all cases, the log-hazard curve shows approximately linear growth in
the late-life regime, confirming that Gompertz emergence is a structural
property of the desynchronization feedback. The exponent $x = 3.5$ was selected
because it produces a biologically plausible lifespan distribution with a clear
plateau before rapid late-life decline; exponents in the range $[2.5, 4.5]$
produce qualitatively identical results.

\textbf{Conclusion:} Aging acceleration in this model emerges from system
architecture — convex cost and feedback-amplified noise — not from the choice
of $x$.

% ─────────────────────────────────────────
\section{Scaling Law Under Non-Mean-Field Topologies}
\label{app:topology}

\subsection{Motivation}

The scaling law $(\Delta\sigma/\sigma)_\text{local}/(\Delta\sigma/\sigma)_\text{global} = 1/n$
was derived under the mean-field assumption: all oscillators are coupled
equally to all others. Real biological regulatory networks are not fully
connected — they have hierarchical structure, hub nodes (e.g., the
suprachiasmatic nucleus as master clock), and modular organization. This
appendix examines how topology affects the $1/n$ scaling.

\subsection{Three topological regimes}

We consider three canonical topologies:

\textbf{(i) Mean-field (all-to-all):} Every oscillator couples to every other
with equal weight $K/n$. This is the baseline case proven in Appendix~A.

\textbf{(ii) Hub-and-spoke:} One master oscillator (hub) couples to all others;
peripheral oscillators do not couple to each other. This approximates the
biological circadian hierarchy, where the suprachiasmatic nucleus entrains
peripheral clocks.

\textbf{(iii) Modular:} Oscillators are partitioned into $m$ modules of size
$n/m$, with strong within-module coupling and weak between-module coupling.
This approximates the organization of major physiological systems (cardiovascular,
metabolic, immune, etc.) that couple loosely to each other.

\subsection{Effect on local intervention efficacy}

The key question is: does a local repair of one oscillator produce a $1/n$
effect, or does topology create amplification or suppression?

\textbf{Hub-and-spoke:} Repairing the hub oscillator affects all $n$ spokes
through the coupling term — the effect is $O(1)$, not $O(1/n)$. Repairing
a single spoke affects only that spoke's contribution to the order parameter,
giving $O(1/n)$ as in the mean-field case. This creates a \emph{heterogeneous}
intervention landscape: hub repair scales as a global intervention, spoke repair
scales locally.

\textbf{Modular topology:} Repairing one oscillator within a module of size
$n/m$ produces a local effect of order $m/n$ — larger than the mean-field
$1/n$ by a factor of $m$. However, this improvement is \emph{confined to the
module}: since between-module coupling is weak, the repaired module does not
pull the rest of the system toward synchrony. The global desynchronization
$\sigma$ is dominated by between-module variance, which the local repair
does not address.

\subsection{General bound}

Across all topologies, the following bound holds:

\begin{equation}
\frac{(\Delta\sigma/\sigma)_\text{local}}{(\Delta\sigma/\sigma)_\text{global}}
\leq \frac{1}{n_\text{eff}}
\end{equation}

where $n_\text{eff}$ is the \emph{effective number of independent regulatory
units} — the number of oscillators that contribute independently to total system
variance. For fully connected networks, $n_\text{eff} = n$. For modular
networks, $n_\text{eff} \approx m$ (number of modules), since within-module
oscillators are highly correlated and contribute as a single unit.

The important consequence: in modular networks, the $1/n$ scaling becomes
$1/m$ — which is \emph{more pessimistic} than mean-field for the same total
$n$, because the effective dimensionality is lower. A human with $n=20$
oscillators organized into $m=6$ physiological modules would show local
intervention efficacy of at most $1/6 \approx 17\%$ per module — but since
each module contains multiple oscillators, repairing one oscillator within a
module gives at most $1/n$ system-wide improvement.

\subsection{Hub repair as global intervention}

The hub-and-spoke analysis reveals a topological exception: interventions
targeting high-centrality nodes (master oscillators, hub genes, master
regulators) can achieve near-global efficacy from a single targeted repair.
This is consistent with the empirical observation that interventions targeting
the circadian master clock (light therapy, chronotherapy) show broader
physiological effects than most molecular interventions.

This does not contradict the main scaling argument. It refines it: the $1/n$
bound applies to \emph{structurally peripheral} components. Interventions
targeting \emph{structural hubs} can exceed this bound — but the number of true
hubs in any biological network is small by definition.

\subsection{Conclusion}

The $1/n$ scaling law is robust across topologies in the following sense: for
interventions targeting non-hub, non-master components (which comprise the vast
majority of molecular targets in current anti-aging research), the system-wide
efficacy scales inversely with the effective number of independent regulatory
units. Topology can create exceptions for high-centrality targets, but these
exceptions are rare by network structure.

This analysis strengthens rather than weakens the main conclusion: the
translational gap between mice and humans is not merely a consequence of
total oscillator count, but of the greater modular complexity of human
regulatory architecture — a complexity that makes hub-targeting strategies
potentially more valuable than pathway-targeting strategies in human aging
research.

% ─────────────────────────────────────────
\section{Directional Test of Desynchronization Predictions in Population Data}
\label{app:nhanes}

\subsection{Dataset}

We analyzed six physiological markers — fasting glucose, total cholesterol,
high-sensitivity CRP, triglycerides, and systolic and diastolic blood pressure
— in NHANES 2017--2018 participants aged 20--79 ($n = 2{,}673$, after exclusion
of participants with any missing marker value). Participants were divided into
six decade-wide age bins (Table~\ref{tab:nhanes}).

\begin{table}[h]
\centering
\caption{Sample sizes per age bin, NHANES 2017--2018.}
\label{tab:nhanes}
\begin{tabular}{lc}
\toprule
Age bin & $n$ \\
\midrule
20--29 & — \\
30--39 & — \\
40--49 & — \\
50--59 & — \\
60--69 & — \\
70--79 & — \\
\hline
Total & 2,673 \\
\bottomrule
\end{tabular}
\end{table}

\subsection{PCA and effective dimensionality}

Principal component analysis on z-scored markers yields cumulative explained
variance:
\begin{equation}
[0.300,\ 0.486,\ 0.662,\ 0.811,\ \mathbf{0.915},\ 1.000]
\end{equation}

The number of components explaining $\geq 90\%$ of variance is $n_\text{eff} = 5$.

This result has an important interpretation: six metabolic markers, despite
representing nominally distinct physiological systems, share substantial
covariance and are effectively spanned by 5 independent dimensions. However,
this panel captures primarily metabolic regulation. A broader panel including
circadian rhythm markers (melatonin, cortisol diurnal profile), immune
parameters, and endocrine axes would be expected to yield higher $n_\text{eff}$,
as these systems are regulated by partially independent mechanisms.

The finding that $n_\text{eff} = 5$ for this metabolic panel does not imply
that human regulatory complexity equals that of mice — it implies that this
particular marker set captures one regulatory module rather than the full
hierarchy. This is consistent with the model's framework: $n$ should be
interpreted as the number of independent regulatory modules relevant to the
intervention target, not total organism complexity.

\subsection{Results}

Within-system variability (mean standard deviation across z-scored markers)
by age decade:
\begin{equation}
[25.6,\ 32.8,\ 37.8,\ \mathbf{48.4},\ 30.7,\ 39.6]
\end{equation}

The general trend is increasing ($R^2 = 0.77$, $p = 0.021$; Figure~\ref{fig:d1}),
consistent with model prediction of rising $\sigma$ with age. The non-monotonic
drop at 60--69 likely reflects survival selection: the least healthy individuals
in this cohort have already died, leaving a relatively healthier survivor
population.

Mean absolute cross-system correlation by age decade:
\begin{equation}
[0.159,\ 0.207,\ 0.157,\ 0.110,\ 0.115,\ 0.117]
\end{equation}

A declining trend is present from the fourth decade onward ($R^2 = 0.56$,
$p = 0.087$; Figure~\ref{fig:d2}), though not reaching conventional significance.
The correlation trend does not reach statistical significance ($p = 0.087$),
likely reflecting the cross-sectional design and limited marker set rather than
absence of the predicted signal.
The elevated value in the 30--39 bin is visible in Figure~D2 and is discussed
below.

\subsection{Interpretation}

These findings do not confirm the synchronization framework — the cross-sectional
design, limited marker selection, and non-monotonic patterns prevent strong
inference. They constitute a \emph{directional test}: the model predicts that
$\sigma$ should increase and cross-system correlation should decrease with age.
Both signals are present in the expected direction.

The $n_\text{eff} = 5$ result from PCA provides an empirical anchor for the
model parameter. For interventions targeting this metabolic regulatory module
specifically, the predicted local efficacy at $f = 0.5$ is $(2f-f^2)/n_\text{eff}
= 0.75/5 = 15\%$ — matching the mouse parameter in the main model. This
reinforces the point that what matters is not total organism complexity but the
effective dimensionality of the \emph{relevant} regulatory module.

The elevated cross-correlation in the 30--39 bin may reflect cohort composition,
medication use patterns, or cross-sectional design effects. Circadian rhythm
amplitude and precision are known to decline with age \citep{monk2005},
consistent with declining cross-system correlation as a marker of regulatory
decoupling.

A longitudinal dataset with repeated individual measurements (e.g.,
\citealt{pyrkov2021}) with direct trajectory fitting would provide a
substantially stronger test of the desynchronization hypothesis.

\newpage
% ─────────────────────────────────────────
\bibliographystyle{plainnat}

\begin{thebibliography}{99}

\bibitem{kuramoto1975}
Kuramoto, Y. (1975).
Self-entrainment of a population of coupled non-linear oscillators.
In \textit{International Symposium on Mathematical Problems in Theoretical Physics},
Lecture Notes in Physics, Vol. 39, pp. 420--422. Springer, Berlin.

\bibitem{pyrkov2021}
Pyrkov, T.V., Avchaciov, K., Tarkhov, A.E., Menshikov, L.I., Gudkov, A.V., \&
Fedichev, P.O. (2021).
Longitudinal analysis of blood markers reveals progressive loss of resilience and
predicts human lifespan limit.
\textit{Nature Communications}, 12, 2765.

\bibitem{strogatz2000}
Strogatz, S.H. (2000).
From Kuramoto to Crawford: exploring the onset of synchronization in populations
of coupled oscillators.
\textit{Physica D: Nonlinear Phenomena}, 143(1--4), 1--20.

\bibitem{barzilai2018}
Barzilai, N., Crandall, J.P., Kritchevsky, S.B., \& Espeland, M.A. (2016).
Metformin as a tool to target aging.
\textit{Cell Metabolism}, 23(6), 1060--1065.

\bibitem{rakshit2023}
Rakshit, K., et al. (2023).
Zeitgeber-based interventions and biological aging: a systematic review of
exercise timing, sleep regularity, and intermittent fasting in animal and
human models.
\textit{Ageing Research Reviews}, 85, 101850.

\bibitem{monk2005}
Monk, T.H. (2005).
Aging human circadian rhythms: conventional wisdom may not always be right.
\textit{Journal of Biological Rhythms}, 20(4), 366--374.

\bibitem{nhanes2018}
Centers for Disease Control and Prevention. (2020).
National Health and Nutrition Examination Survey Data 2017--2018.
U.S. Department of Health and Human Services.
\url{https://www.cdc.gov/nchs/nhanes/}

\bibitem{lopez2023hallmarks}
L\'{o}pez-Ot\'{i}n, C., Blasco, M.A., Partridge, L., Serrano, M., \&
Kroemer, G. (2023).
Hallmarks of aging: An expanding universe.
\textit{Cell}, 186(2), 243--278.

\end{thebibliography}

\end{document}
